\section{Experimental Design}
\label{sec:experiments}

\subsection{Linear Scaling Rule}

All learning rates are derived from the Linear Scaling Rule of
Goyal~et~al.~\citep{goyal2017accurate}:

\begin{equation}
  \eta = \eta_\text{base} \times \frac{B_\text{total}}{B_\text{ref}},
  \label{eq:lsr}
\end{equation}

\noindent where $B_\text{total}$ is the total effective batch size
(per-GPU batch $\times$ number of GPUs), $B_\text{ref}$ is the
reference batch size, and $\eta_\text{base}$ is the corresponding
reference learning rate.  Reference values are:

\begin{itemize}[leftmargin=1.5em]
  \item \textbf{PicoDet-S:} $\eta_\text{base}=0.08$, $B_\text{ref}=32$
    (single-GPU baseline).
  \item \textbf{YOLOv3:} $\eta_\text{base}=0.00125$, $B_\text{ref}=16$
    (single-GPU baseline, per-GPU batch=16).
\end{itemize}

A linear warm-up over 300 steps is applied at the start of all runs,
followed by cosine decay.

\subsection{Experiment Matrix}

\subsubsection{PicoDet-S (L-series, full dataset)}

Four runs vary GPU count, per-GPU batch size, and training duration
while keeping the dataset fixed at 5,828 training images.

\begin{table}[H]
\centering
\caption{PicoDet-S experiment configurations (L-series).}
\label{tab:picodet_config}
\resizebox{\textwidth}{!}{%
\begin{tabular}{lccccccc}
\toprule
\textbf{Run} & \textbf{GPUs} & \textbf{BS/GPU} & \textbf{Total BS}
             & \textbf{LR} & \textbf{Epochs}
             & \textbf{Steps/Epoch} & \textbf{Description} \\
\midrule
L1 & 1 & 32 & 32  & 0.08 & 300 & 182 & Single-GPU baseline \\
L2 & 2 & 32 & 64  & 0.16 & 300 &  91 & Dual-GPU standard \\
L3 & 2 & 48 & 96  & 0.24 & 300 &  61 & Large batch \\
L4 & 2 & 32 & 64  & 0.16 & 600 &  91 & Extended training \\
\bottomrule
\end{tabular}}
\end{table}

\subsubsection{YOLOv3 (Y-series, variable dataset)}

Four runs vary both GPU count and training set size in a $2\!\times\!2$
factorial design.  All runs use 80 epochs and per-GPU batch size of 16.

\begin{table}[H]
\centering
\caption{YOLOv3 experiment configurations (Y-series).}
\label{tab:yolov3_config}
\resizebox{\textwidth}{!}{%
\begin{tabular}{lcccccc}
\toprule
\textbf{Run} & \textbf{GPUs} & \textbf{Total BS}
             & \textbf{LR} & \textbf{Train Images} & \textbf{Epochs}
             & \textbf{Description} \\
\midrule
Y1 & 1 & 16  & 0.00125 & 1{,}942 & 80 & 1/3 data, single GPU \\
Y2 & 2 & 32  & 0.00250 & 1{,}942 & 80 & 1/3 data, dual GPU \\
Y3 & 1 & 16  & 0.00125 & 5{,}828 & 80 & Full data, single GPU \\
Y4 & 2 & 32  & 0.00250 & 5{,}828 & 80 & Full data, dual GPU \\
\bottomrule
\end{tabular}}
\end{table}

\subsection{Infrastructure}

All experiments run on a server equipped with two NVIDIA Tesla T4
(15\,GB VRAM each).  Multi-GPU communication uses the GLOO backend
over TCP (NCCL was unavailable due to the absence of InfiniBand),
which adds negligible overhead for our model sizes.
The measured dual-card speed-up is 1.97--2.22$\times$, close to the
theoretical $2\times$ limit.

Training is orchestrated by two chain-execution Python scripts
(\texttt{run\_lightweight.py} for L-series,
\texttt{run\_yolov3.py} for Y-series) that launch runs
sequentially, log all metrics to per-run directories under
\texttt{output/}, and write consolidated CSV summaries.

\subsection{Evaluation Protocol}

Model accuracy is measured as mAP@0.5 (mean Average Precision at
IoU threshold 0.5) using the COCO-style integral computation
provided by PaddleDetection's evaluation pipeline.
Evaluation is performed every 10 epochs on the shared 1,458-image
validation set; \texttt{best\_model.pdparams} is updated whenever
a new peak is reached.
FPS is measured on the same Tesla T4 server using the framework's
built-in profiler with batch size~1.
